\documentclass[12pt]{article} 
\usepackage{imakeidx}
\usepackage{graphicx}
\usepackage{soul}
\usepackage{hyperref}

% Keep hyperlink text black and remove PDF link-border boxes.
\hypersetup{
  colorlinks=true,
  linkcolor=black,
  urlcolor=black,
  citecolor=black
}

% Make link underlines blue while keeping the text black.
\setulcolor{blue}

\title{An historical and mathematical analysis of Bitonic sort}
\author{Pietro Marini, Michele Riolfi, Javidan Tahmazli}
\date{November 2026}

% To add a heading to the index, place its matching index command immediately after the heading and its unique \label. Use \indexsection{KEY}{LABEL}{TITLE}

% The section key describes the order among entries at the same level.
% This must be consistent with the way the sections appear in the rest of the document.

% For a section; write as the KEY a three-digit number that sets its order (001, 002, 003, ...). 

% Use \indexsubsection{SECTION-KEY}{SECTION-LABEL}{SECTION-TITLE}
% {SUBSECTION-KEY}{SUBSECTION-LABEL}{SUBSECTION-TITLE} for a subsection, and \indexsubsubsection with those same section and subsection details followed by  the subsubsection's key, label, and title. 

% Number subsection keys in the desired order within their section, and subsubsection keys within their subsection. Keep each key consistent between the parent heading and its children; make every label unique. The title arguments should match the heading text, since they are what readers see in the index.

% Use one column and the custom style file.
\makeindex[columns=1, options= -s indexstyle.ist]

% Override article's default two-column index layout.
% Because we are generating the .ind file and then inserting it, \makeindex[column=1, ...] doesn't work by itself.
% In fact, it is redundant but kept for clarity.
\makeatletter
\renewenvironment{theindex}
  {%
    \clearpage
    \section*{\indexname}%
    \@mkboth{\MakeUppercase\indexname}{\MakeUppercase\indexname}%
    \thispagestyle{plain}%
    \parindent\z@
    \parskip\z@ plus .3\p@\relax
    \let\item\@idxitem
    \normalsize
  }
  {\clearpage}
\makeatother

% Keep linked page numbers together on one line.
\let\OriginalHyperpage\hyperpage
\renewcommand*{\hyperpage}[1]{\mbox{\OriginalHyperpage{#1}}}

\setcounter{secnumdepth}{0} % Suppress heading numbers
\renewcommand{\indexname}{Index}

% Add a section entry.
% #1 is its numeric sort key; #2 its label; #3 its displayed title.
\newcommand{\indexsection}[3]{%
  \index{#1@\hyperref[#2]{\ul{#3}}}%
}

% Add a subsection beneath its parent section.
% Arguments are: section key, label, title; subsection key, label, title.
\newcommand{\indexsubsection}[6]{%
  \index{#1@\hyperref[#2]{\ul{#3}}!%
    #4@\hyperref[#5]{\ul{#6}}}%
}

% Add a subsubsection beneath its parent subsection.
% Arguments are: section key, label, title;
% subsection key, label, title; subsubsection key, label, title.
\newcommand{\indexsubsubsection}[9]{%
  \index{#1@\hyperref[#2]{\ul{#3}}!%
    #4@\hyperref[#5]{\ul{#6}}!%
    #7@\hyperref[#8]{\ul{#9}}}%
}

\begin{document}

\maketitle
\clearpage % Force new page

% Insert index.
% \printindex by itself does not work because the package only allows for the index to be at the end.
% This workaround solution makes it so MakeIndex is ran after the LaTex pass and then the generated .ind file is output in page 2.
% The two \clearpage are there just so the index has a dedicated page.
\clearpage
\IfFileExists{\jobname.ind}{\input{\jobname.ind}}{}
\clearpage

\section{History}\phantomsection\label{sec:history}
\indexsection{001}{sec:history}{History}

\subsection{Original paper}\phantomsection\label{subsec:original-paper}
\indexsubsection{001}{sec:history}{History}
                 {001}{subsec:original-paper}{Original paper}

The \href{https://dl.acm.org/doi/epdf/10.1145/1468075.1468121}
{\ul{original paper}} is hosted by the Association for Computing Machinery
digital library.

\subsection{Related context, terminology and developments}
\phantomsection\label{subsec:related-context}
\indexsubsection{001}{sec:history}{History}
                 {002}{subsec:related-context}
                 {Related context, terminology and developments}

\subsection{Necessity and details}
\phantomsection\label{subsec:necessity}
\indexsubsection{001}{sec:history}{History}
                 {003}{subsec:necessity}{Necessity and details}

\subsection{Historical reconstruction and implementation}
\phantomsection\label{subsec:historical-reconstruction}
\indexsubsection{001}{sec:history}{History}
                 {004}{subsec:historical-reconstruction}
                 {Historical reconstruction and implementation}

\section{Validation and analysis}\phantomsection\label{sec:validation}
\indexsection{002}{sec:validation}{Validation and analysis}

\subsection{Theoretical analysis}
\phantomsection\label{subsec:theoretical-analysis}
\indexsubsection{002}{sec:validation}{Validation and analysis}
                 {001}{subsec:theoretical-analysis}{Theoretical analysis}

\subsection{Practical analysis}
\phantomsection\label{subsec:practical-analysis}
\indexsubsection{002}{sec:validation}{Validation and analysis}
                 {002}{subsec:practical-analysis}{Practical analysis}

\subsection{Complexity and performance analysis}
\phantomsection\label{subsec:complexity}
\indexsubsection{002}{sec:validation}{Validation and analysis}
                 {003}{subsec:complexity}{Complexity and performance analysis}

\subsubsection{Chosen experiments and related info}
\phantomsection\label{subsubsec:experiments}
\indexsubsubsection{002}{sec:validation}{Validation and analysis}
                    {003}{subsec:complexity}{Complexity and performance analysis}
                    {001}{subsubsec:experiments}{Chosen experiments and related info}

\section{Reality check, evolution and legacy}
\phantomsection\label{sec:reality-check}
\indexsection{003}{sec:reality-check}{Reality check, evolution and legacy}

\end{document}
